\gostunnumberedsection{Введение}{Введение}
\label{sec:intro}

Дисбаланс классов является одной из наиболее распространённых и трудноустранимых особенностей данных в прикладных задачах машинного обучения. В типичных постановках доля миноритарного класса (событие, дефолт, мошенничество, отклик на предложение) на порядки и более раз меньше доли мажоритарного класса. При таком соотношении классов модели, обученные минимизацией стандартных функций потерь, склонны к вырожденному решению, предсказывающему доминирующий класс, что приводит к нулевой или неприемлемо низкой полноте по редкому классу при формально высокой доле верных ответов.

Актуальность темы обусловлена масштабом внедрения моделей машинного обучения в финансовом секторе. Кредитный скоринг, антифрод-системы, маркетинговые кампании с таргетированием, раннее выявление оттока клиентов и мониторинг операционных рисков опираются на классификацию с сильным перекосом частот. Ошибка второго рода (пропуск редкого события) и ошибка первого рода (ложная тревога) имеют различную экономическую стоимость, поэтому универсальные методы, не учитывающие предметную область, часто не достигают целевых показателей качества в промышленной эксплуатации.

Состояние исследований в области дисбаланса классов характеризуется высокой степенью разработанности отдельных групп методов и одновременным отсутствием единого подхода, гарантирующего успех на произвольных данных. На уровне данных широко применяются методы перевыборки (oversampling, undersampling, гибридные схемы), в том числе алгоритм SMOTE и его многочисленные модификации. На уровне алгоритмов используются взвешивание классов, cost-sensitive learning, специализированные функции потерь (в частности, focal loss), а также ансамблевые и каскадные архитектуры. Для оценки качества предложены метрики, ориентированные на редкий класс: precision-recall AUC, balanced accuracy, F-beta при различных beta, а также кривые стоимости ошибок при заданных бизнес-ограничениях.

Промышленные внедрения показывают, что перенос методов между предметными областями без адаптации часто даёт нестабильный эффект. В прикладных доменах с табличными данными (скoring, antifraud, таргетинг) наблюдаются дрейф распределений, сезонность, смещение выборки (selection bias), отложенная разметка и высокая стоимость экспертной проверки. Классический pipeline «перевыборка --- градиентный бустинг» в таких условиях не исчерпывает задачу; требуются подходы, явно учитывающие предметную специфику: бизнес-метрики, ограничения на latency и interpretability, мониторинг и воспроизводимость экспериментов.

Степень разработанности проблемы. Фундаментальные работы по learning from imbalanced data систематизированы в монографиях и обзорах He и Garcia, Krawczyk, а также в специализированных журналах и материалах конференций IEEE ICDM, KDD, NeurIPS. Метод SMOTE (Chawla et al.) и его развитие (Borderline-SMOTE, ADASYN, SMOTE-ENN) остаются базовым инструментом инженерной практики. Для табличных данных доминируют градиентные бустинги (XGBoost, LightGBM, CatBoost) с настройкой весов и порогов решения. В последние годы активно исследуются deep learning подходы, metric learning и генеративные модели для синтеза миноритарных примеров, а также методы на основе partial labels и weak supervision.

Обзор ключевых работ показывает, что предметно-ориентированные элементы чаще появляются в виде эвристик (ручной подбор порога по validation, фиксированные cost matrices) и реже --- в виде формализованных методов, интегрирующих domain knowledge в процесс обучения и калибровки. Для финансового сектора описаны кейсы fraud detection и credit scoring, однако открытыми остаются вопросы переносимости решений между продуктами, устойчивости к covariate shift и согласования offline-метрик с online KPI.

Цель исследования --- разработать предметно-ориентированные подходы к решению задач классификации с дисбалансом классов и экспериментально подтвердить их эффективность на данных, приближённых к промышленным постановкам партнёра.

Для достижения цели поставлены следующие задачи:
\begin{enumerate}
    \item провести систематический обзор и сравнительный анализ методов обработки дисбаланса классов на уровне данных, алгоритмов и метрик качества;
    \item выявить специфические ограничения и требования прикладных задач классификации с дисбалансом классов;
    \item сформулировать требования к предметно-ориентированному подходу и разработать метод (алгоритм) с теоретическим обоснованием границ применимости;
    \item реализовать программный комплекс экспериментов и протокол сравнения с базовыми и современными методами (state-of-the-art);
    \item провести экспериментальную оценку предложенного подхода, интерпретировать результаты и оценить воспроизводимость;
    \item сформулировать рекомендации по внедрению результатов в продуктовые процессы.
\end{enumerate}

Объект исследования --- процессы построения и эксплуатации моделей машинного обучения для задач бинарной и многоклассовой классификации при дисбалансе классов.

Предмет исследования --- методы, алгоритмы и программные средства предметно-ориентированной обработки дисбаланса классов, включая выбор метрик, калибровку порогов и интеграцию бизнес-ограничений в цикл обучения.

Научная новизна работы заключается в разработке предметно-ориентированного подхода, объединяющего формализацию бизнес-стоимости ошибок классификации, адаптивную стратегию балансировки и протокол экспериментальной валидации, ориентированный на промышленные KPI (в том числе precision at fixed recall и экономический gain при заданных операционных ограничениях). Отличие от известных методов состоит не в замене отдельного алгоритма перевыборки, а в системной связке этапов: от постановки задачи и выбора метрик до мониторинга стабильности решения во времени.

Теоретическая значимость состоит в уточнении критериев сопоставимости методов при дисбалансе классов и в описании условий, при которых комбинация cost-sensitive обучения и предметно-ориентированной калибровки порога доминирует стандартные baseline-схемы.

Практическая значимость определяется возможностью использования результатов в прикладных задачах скоринга, antifraud и персонализации, а также переносом программной реализации и протокола экспериментов в корпоративные MLOps-процессы.

Методология и методы исследования. В работе используются методы системного анализа, сравнительного эксперимента, statistical learning theory и design of experiments. Теоретический анализ опирается на обзор литературы и классификацию методов. Эмпирическая часть включает вычислительные эксперименты на открытых и промышленных (обезличенных) наборах данных, сравнение с baseline (Logistic Regression, Random Forest, Gradient Boosting), методами перевыборки и современными cost-sensitive решениями. Качество оценивается по метрикам PR-AUC, ROC-AUC, F1, recall at precision, а также по предметным экономическим показателям. Программная реализация выполняется на языке Python с использованием библиотек scikit-learn, imbalanced-learn, PyTorch (при необходимости), MLflow для фиксации экспериментов.

Структура работы. Введение обосновывает актуальность, формулирует цель, задачи, объект, предмет, новизну и методологию. В первой главе приведён обзор методов работы с дисбалансом классов и предметной специфики прикладных задач, сформулированы требования к предлагаемому решению. Во второй главе описаны концепция, формальная постановка и программная реализация предметно-ориентированного подхода, а также план экспериментов. В третьей главе представлены результаты экспериментальной проверки, сравнение с базовыми методами, обсуждение и оценка воспроизводимости. В заключении сформулированы выводы по главам, оценка достижения цели и направления дальнейших исследований. Список использованных источников содержит не менее 60 наименований. В приложении приведена ссылка на репозиторий с исходным кодом и вспомогательные материалы.
